\documentclass[11pt,a4paper]{article}
\usepackage[T1]{fontenc}
\usepackage[utf8]{inputenc}
\usepackage{lmodern,microtype}
\usepackage[margin=28mm]{geometry}
\usepackage{amsmath,amssymb,amsthm,mathtools}
\usepackage{booktabs,tabularx,array,enumitem}
\usepackage{xcolor,xurl}
\usepackage[colorlinks=true,linkcolor=blue!45!black,citecolor=blue!45!black,urlcolor=blue!45!black]{hyperref}
\usepackage{bookmark}
\setlength{\parindent}{1.1em}
\setlength{\parskip}{2pt}
\setlength{\emergencystretch}{2em}
\setlist[enumerate]{itemsep=3pt,topsep=5pt}
\setlist[itemize]{itemsep=2pt,topsep=4pt}
\numberwithin{equation}{section}
\newtheorem{proposition}{Proposition}[section]
\newtheorem{lemma}[proposition]{Lemma}
\newtheorem{theorem}[proposition]{Theorem}
\newtheorem{corollary}[proposition]{Corollary}
\theoremstyle{definition}
\newtheorem{definition}[proposition]{Definition}
\newtheorem{notation}[proposition]{Notation}
\theoremstyle{remark}
\newtheorem{remark}[proposition]{Remark}
\newtheorem{example}[proposition]{Example}
\newcommand{\Z}{\mathbf Z}
\newcommand{\Q}{\mathbf Q}
\newcommand{\R}{\mathbf R}
\newcommand{\GL}{\operatorname{GL}}
\newcommand{\St}{\operatorname{St}}
\newcommand{\Vor}{\operatorname{Vor}}
\newcommand{\Sym}{\operatorname{Sym}}
\newcommand{\sgn}{\operatorname{sgn}}
\newcommand{\rank}{\operatorname{rank}}
\newcommand{\im}{\operatorname{im}}
\newcommand{\coker}{\operatorname{coker}}
\newcommand{\Span}{\operatorname{span}}
\newcommand{\tr}{\operatorname{tr}}
\newcommand{\BQ}{BQ\Z}
\title{\vspace{-1.2em}Verifying the computation of \(K_8(\Z)\):\\
finite certificates and admitted inputs}
\author{Project synthesis}
\date{28 October 2026}
\begin{document}
\maketitle
\vspace{-1.3em}
\begin{abstract}
We describe a local verification of the rank-eight Voronoi calculation
used in the proof that \(K_8(\Z)=0\).
The finite task is to exhibit an orientation-reversing integral symmetry
for each of 106 published vector configurations.
We specify a finite search, an independent exact checker, and the
mathematical implication of a successful run.
Completeness of the cell list, lower-rank computations, and the
comparison and arithmetic theorems remain explicit inputs.
The result is a reproducible verification relative to these inputs;
the document does not assume that a successful run has already occurred.
\end{abstract}

\section{The question and the strategy}\label{sec:intro}

The difficulty in computing \(K_8(\Z)\) is the possible torsion.
Calculations with rational coefficients alone cannot exclude it.
Dutour Sikirić, Elbaz-Vincent, Kupers and Martinet proved that
\(K_8(\Z)=0\) by combining low-dimensional Voronoi cells,
earlier arithmetic-group homology calculations, and arithmetic
information at small primes \cite[Theorem~6.4]{DEKM}.
Our purpose is to make a specific finite ingredient of this proof
independently checkable on an ordinary computer.

Throughout the report let
\[
 R=\Z[1/7!]=\Z[1/210].
\]
The Voronoi comparison turns the required rank-eight calculation into
\(H_8(\Vor_8;R)=0\). In this case there is a particularly economical
certificate: every generator in chain degree eight disappears because
its stabilizer reverses its orientation. The task therefore consists
of finding and checking finitely many integral matrices. No boundary
matrix or Smith reduction is needed for this step.

There are two distinct questions about the input. First, does the
published list contain a representative of every relevant cell?
Second, does each listed representative possess the asserted symmetry?
We admit the published classification for the first question and
compute witnesses for the second. Keeping these questions separate
is essential: checking a finite file cannot by itself prove that
no orbit was omitted from its construction.

Section~\ref{sec:data} identifies the geometric input and fixes the
degree convention. Proposition~\ref{prop:orientation} gives the exact
orientation test, while Proposition~\ref{prop:search} proves that the
specified finite search finds a reversing symmetry whenever one exists.
Section~\ref{sec:bridge} records the published inputs that pass from
Voronoi homology to \(K\)-theory.
The final Theorem~\ref{thm:certificate} states precisely what a complete
successful run certifies, and under which premises it recovers
\(K_8(\Z)=0\).

\section{The cells and the admitted data}\label{sec:data}

Let \(L=\Z^N\), and regard lattice vectors as columns. For a positive
definite quadratic form on \(L\), choose one vector from each pair
of opposite minimal vectors. Their rank-one matrices generate a cone
in \(\Sym_N(\R)\). The cones arising in this way form the relevant
faces of the perfect-cone decomposition; passage to positive scaling
gives the projective cells of the Voronoi complex
\cite[\S1.2]{DEKM}.

\begin{notation}\label{not:configuration}
For a finite configuration \(V=(v_1,\ldots,v_s)\) of nonzero integral
vectors, let us write
\[
 A_i=v_i v_i^{\,t},\qquad
 W_V=\Span_{\Q}\{A_1,\ldots,A_s\},\qquad
 \sigma_V=\sum_{i=1}^s\R_{\geq0}A_i.
\]
We write \(c=\dim_{\Q}W_V\). The cone has dimension \(c\), and its
projectivization has dimension \(c-1\).
\end{notation}

The ambient rank \(N\), the number \(s\) of vector pairs, and the cone
dimension \(c\) have different meanings. In particular, the target
\(H_8(\Vor_8;R)\) concerns \(N=8\) and \(c=9\), not \(c=8\).
The program will calculate both the rank of the vectors and the rank
of their rank-one matrices over \(\Q\).

The proposed input is the file
\[
 \text{\small\path{DATA_gap/ListSHV_n8_rnk9}}
\]
in the authors' database \cite{Data}. A preliminary inspection of the
file gives 106 records. Each record contains nine primitive vectors
of \(\Z^8\), spanning \(\Q^8\), whose rank-one matrices are linearly
independent. Thus \(N=8\), \(s=c=9\), and the projective dimension is
eight. These elementary properties will be recomputed by the verifier;
this inspection is not a symmetry calculation.

The file labels require care. The neighboring file with label
\texttt{rnk8}, and the standalone \texttt{cells\_GL8\_dim8.gp}, concern
the preceding cone dimension. One must use the computed rank \(c\)
to identify the homological degree. This also avoids transferring
a dimension label from a table without checking its convention.

\begin{definition}\label{def:admission}
Let us define the \emph{classification input} to be the assertion that
the chosen list consists of realizable Voronoi cells and contains a
representative of every \(\GL_8(\Z)\)-orbit of projective eight-cells
meeting the positive definite locus.
\end{definition}

Version one admits Definition~\ref{def:admission} from the published
classification \cite[\S\S3--4]{DEKM} and the accompanying database
\cite{Data}. This includes the identification of the retained file
with the published classification. It does not enumerate all perfect forms,
test realizability by shortest-vector searches, or independently prove
that every orbit occurs. Pairwise inequivalence of the representatives
is unnecessary for the vanishing argument: coverage is what matters.

The input must nevertheless have a reproducible identity.
The implementation will retain the original file, its repository
revision and checksum, and a deterministic conversion to column
vectors. It will check integrality, primitivity, absence of zero or
duplicate rays, full ambient rank, and cone dimension.
The checksum fixes which bytes were checked; the classification input
supplies the mathematical completeness assertion.

\section{An exact orientation certificate}\label{sec:orientation}

Let \(V\) span \(\Q^N\), and suppose its vectors represent distinct
primitive lattice lines. An integral symmetry is a matrix
\(g\in\GL_N(\Z)\) for which
\begin{equation}\label{eq:signed}
 g v_i=\epsilon_i v_{p(i)}
 \qquad(1\leq i\leq s),
 \qquad \epsilon_i\in\{1,-1\},\quad p\in\mathfrak S_s.
\end{equation}
It acts on \(W_V\) by \(T_g(A)=gAg^t\).
This convention corresponds to the right action
\(A\mapsto\gamma^tA\gamma\) in the source by taking \(\gamma=g^t\).

The sign of \(\det(g)\) on the lattice is not the orientation sign
of the cell. The latter concerns the action on symmetric matrices.

\begin{definition}\label{def:certificate}
Let us define an \emph{orientation certificate} for \(V\) to be an
integer matrix \(g\), a permutation \(p\), and signs \(\epsilon_i\)
satisfying \eqref{eq:signed}, together with the conditions
\[
 \det(g)\in\{1,-1\},
 \qquad
 \det\!\left(T_g|_{W_V}\right)<0.
\]
A saved certificate also records the source configuration and the
exact coordinate calculation of the second determinant.
\end{definition}

To compute this determinant, choose the first linearly independent
\(A_i\)'s in a fixed ordering, obtaining a basis
\(B_1,\ldots,B_c\) of \(W_V\). Exact elimination gives the matrix
\(M_g\) determined by
\begin{equation}\label{eq:action}
 gB_k g^t=\sum_{\ell=1}^{c}(M_g)_{\ell k}B_\ell.
\end{equation}
One may flatten a symmetric matrix into its upper-triangular entries;
this is simply a coordinate isomorphism and requires no approximation.
The checker verifies the equalities in \eqref{eq:action} and computes
\(\det M_g\) over \(\Q\).

\begin{proposition}\label{prop:orientation}
An orientation certificate reverses the orientation of the projective
cell associated with \(V\). If the matrices \(A_1,\ldots,A_s\) are
linearly independent, its orientation sign is \(\sgn(p)\).
\end{proposition}

\begin{proof}
Put \(C=\sum_i A_i\). Equation~\eqref{eq:signed} gives
\(gCg^t=C\). The matrix \(C\) lies in the relative interior of
\(\sigma_V\), and the tangent space of its projectivization at
\([C]\) is \(W_V\otimes\R\) modulo the fixed line \(\R C\).
The action on that line is the identity. Its removal therefore does
not change the determinant sign.

If the \(A_i\)'s are independent, they form a basis and
\(T_g(A_i)=A_{p(i)}\). Thus \(M_g\) is the permutation matrix of \(p\).
The signs \(\epsilon_i\) disappear on passing to rank-one matrices.
\end{proof}

For the proposed 106 configurations, the independence check makes the
second statement applicable: an odd permutation of the nine rays
suffices. Computing \eqref{eq:action} as well provides a direct check
of the degree and orientation conventions. For a more general
configuration with relations between its rank-one generators, the
action on their span remains the correct test.

\begin{example}\label{ex:orientation}
Take \(V=(e_1,e_2)\) in \(\Z^2\). The matrix exchanging \(e_1\) and
\(e_2\) exchanges the two generators of the cone of nonnegative
diagonal matrices, and reverses its projective interval.
By contrast, \(g=\operatorname{diag}(-1,1)\) has lattice determinant
\(-1\) but fixes both rank-one generators. It preserves the cell's
orientation. This example separates the two determinants.
\end{example}

\section{The finite search and the independent checker}\label{sec:algorithm}

We now specify a complete finite search. It supplies a precise
baseline algorithm; an implementation may use a faster symmetry
routine to produce the same kind of witnesses.

For each configuration let \(C=\sum_i v_i v_i^t\) and \(Q=C^{-1}\).
Full ambient span makes \(C\) positive definite.
Any solution of \eqref{eq:signed} satisfies
\begin{equation}\label{eq:isometry}
 gCg^t=C,\qquad g^tQg=Q.
\end{equation}
The rational scalar products defined by \(Q\) therefore constrain
the possible images of a basis.

Choose deterministically \(N\) independent vectors
\(u_1,\ldots,u_N\) from \(V\), and put \(U=(u_1\ \cdots\ u_N)\).
Run the following procedure, with exact integer and rational arithmetic.
\begin{enumerate}
\item Choose the images \(w_a\) successively from
      \(\{\pm v_1,\ldots,\pm v_s\}\), without repeating a lattice line.
      Keep a partial choice only if
      \[
       w_a^tQw_b=u_a^tQu_b
      \]
      for every pair of indices already chosen, including \(a=b\).
\item For each surviving complete choice, put
      \(U'=(w_1\ \cdots\ w_N)\) and \(g=U'U^{-1}\).
      Reject it unless all entries of \(g\) are integers and
      \(\det g=\pm1\).
\item Check the image of every vector in \(V\), recovering and
      verifying a bijection \(p\) and the signs in \eqref{eq:signed}.
      Compute the induced action \eqref{eq:action}.
\item Save the first candidate with negative orientation sign.
      If the entire finite search has no such candidate, report that
      fact for this configuration. If a resource limit interrupts
      the search, report it as incomplete.
\end{enumerate}
Using a basis of the rational span does not mean that \(U\) is a
unimodular matrix. The integrality and determinant checks in the
second step are indispensable.

\begin{proposition}\label{prop:search}
Without a resource cutoff, this procedure terminates and finds an
orientation certificate whenever one exists.
\end{proposition}

\begin{proof}
There are at most
\(2^N s!/(s-N)!\) choices of signed images of the independent vectors.
Every integral symmetry gives one of these choices, and
\eqref{eq:isometry} shows that none of its partial choices is removed
by the scalar-product tests. The images of a basis determine the
linear map uniquely, so the second step recovers that symmetry.
The remaining tests are exactly its defining conditions.
\end{proof}

This finite bound is not a practical runtime estimate.
The scalar-product constraints can remove most branches early;
their effectiveness must be measured on the actual data.
For a successful proof by witnesses, the search need not exhaust
the stabilizer: one reversing element per cell is enough.

The independent checker takes the original normalized configurations
and the saved witnesses. For each record it recomputes the ranks,
the unimodularity of \(g\), every equality in \eqref{eq:signed},
and the induced determinant. It does not trust dimensions, signs,
or action matrices merely because the search saved them.
It also verifies that every required input record has a certificate
and that the dataset identity agrees.

A faster search may therefore use group packages or heuristics without
enlarging the mathematical claim accepted by the checker. An incorrect
candidate is rejected. A search timeout leaves an unfinished cell.
Only explicit exact witnesses contribute to the successful coverage
count. The remaining computational trust consists of the checker,
its arithmetic libraries, and their execution; this is not a
formal proof-assistant verification.

\section{The admitted passage to \texorpdfstring{\(K\)}{K}-theory}\label{sec:bridge}

We record the published inputs in the order in which the conclusion
uses them. They are mathematical premises of the local verification,
not theorems that the program proves.

Let \(\St_N\) be the Steinberg module: the top reduced integral homology
of the rational Tits building for \(N\geq2\), with the customary
rank-one convention. The equivariant cell spectral sequence for the
Voronoi decomposition gives
\begin{equation}\label{eq:comparison}
 H_j(\Vor_N;R)\simeq
 H_{j-N+1}\bigl(\GL_N(\Z);\St_N\otimes R\bigr)
\end{equation}
when the finite cell-stabilizer orders are invertible in \(R\)
\cite[Proposition~1.5]{DEKM}. In rank eight their prime divisors
are at most nine, hence lie in \(\{2,3,5,7\}\).
The chosen ring \(R\) therefore suffices.

More explicitly, the contribution of a cell to the degree-zero
stabilizer row is the module of coinvariants of its orientation line.
A reversing stabilizer identifies a generator \(x\) with \(-x\);
since \(2\) is invertible, this module is zero.
Positive stabilizer homology also vanishes over \(R\).
Thus certificates for every projective eight-cell imply
\begin{equation}\label{eq:rank8}
 V_8(\Vor_8;R)=0,\qquad
 H_8(\Vor_8;R)=0,\qquad
 H_1(\GL_8(\Z);\St_8\otimes R)=0.
\end{equation}
The first equality needs only completeness of the cell list and
the checked symmetries; the comparison theorem supplies the last.

Quillen's rank filtration gives the first-quadrant spectral sequence
\begin{equation}\label{eq:rankfiltration}
 E^1_{N,q}=H_q(\GL_N(\Z);\St_N\otimes R)
 \ \Longrightarrow\ H_{N+q}(\BQ;R),
\end{equation}
with the rank-zero starting space a point
\cite[Theorem~3]{Quillen}; see also
\cite[\S6, proof of Proposition~6.1]{DEKM}.
The complete diagonal needed in total degree nine is as follows.
\begin{center}
\small
\begin{tabularx}{\linewidth}{@{}c c X@{}}
\toprule
Rank \(N\)&Degree \(q\)&Input giving \(E^1_{N,q}=0\) over \(R\)\\
\midrule
\(1\)&\(8\)&
\(\GL_1(\Z)=\{1,-1\}\), whose positive homology vanishes after
inverting \(2\).\\
\(2,3,4\)&\(7,6,5\)&
The classical lower-rank results recorded in
\cite[Table~5]{DEKM}.\\
\(5,6,7\)&\(4,3,2\)&
The vanishing of degree-eight Voronoi homology in
\cite[Theorem~4.3]{EVGS}, with the comparison \eqref{eq:comparison}.\\
\(8\)&\(1\)&
The checked orientation certificates, the classification input,
and \eqref{eq:rank8}.\\
\(9\)&\(0\)&
Steinberg coinvariant vanishing
\cite[Theorem~1.3]{LS}; see \cite[Remark~6.2]{DEKM}.\\
\bottomrule
\end{tabularx}
\end{center}
Version one admits the entries outside rank eight.
Ranks above nine cannot contribute because they would require
\(q<0\). Every entry on the diagonal is already zero on the first
page, so there is no differential or extension to determine there.
Consequently,
\begin{equation}\label{eq:H9}
 H_9(\BQ;R)=0.
\end{equation}

By definition \(K_8(\Z)=\pi_9(\BQ)\).
The degree-nine Hurewicz homomorphism
\[
 \pi_9(\BQ)\longrightarrow H_9(\BQ;\Z)
\]
has kernel supported at \(2,3,5\)
\cite[Proposition~6.3]{DEKM}, by the infinite-loop-space
Hurewicz bound of \cite[Theorem~1.5]{Arlettaz}.
It becomes injective after tensoring with \(R\).
Equation~\eqref{eq:H9} therefore implies
\begin{equation}\label{eq:localizedK}
 K_8(\Z)\otimes R=0.
\end{equation}

Finally, the arithmetic inputs remove the four inverted primes.
The group \(K_8(\Z)\) has no \(2\)-primary torsion
\cite[VI, Corollary~9.8]{Weibel}.
For an odd regular prime \(\ell\), the positive even \(K\)-groups of
\(\Z\) have no \(\ell\)-torsion
\cite[VI, Proposition~10.5]{Weibel}; this applies to \(3,5,7\).
These are substantive external theorems. Their role is not replaced
by reducing the cell matrices modulo these primes.

\section{The precise certification claim}\label{sec:claim}

A run should produce a witness for each input record and enough exact
data to recheck it independently: the integral matrix, the signed
permutation, a basis of \(W_V\), and its induced action.
The summary should give the input revision, the number of records
checked, their computed dimensions, any unfinished records, and the
admitted premises. Elapsed time and resource measurements describe
that run; they are not mathematical evidence of completeness.

The logical distinction is between three levels. A checked witness
certifies an exact symmetry of one supplied configuration.
Checking all records certifies the corresponding finite conjunction.
The classification and comparison inputs then give a statement
about the whole Voronoi complex, and the remaining published inputs
give the statement about \(K\)-theory.

\begin{theorem}\label{thm:certificate}
Suppose that the independent checker accepts an orientation certificate
for every record of a list of rank-eight configurations with
projective dimension eight.

For each record, its output certifies an orientation-reversing element
of \(\GL_8(\Z)\) stabilizing the associated cone.
If the classification input of Definition~\ref{def:admission} holds,
then
\[
 H_8(\Vor_8;R)=0.
\]
Under the comparison, lower-rank, rank-filtration, Hurewicz, and
arithmetic inputs specified in Section~\ref{sec:bridge}, it follows
that
\[
 \boxed{K_8(\Z)=0.}
\]
\end{theorem}

\begin{proof}
The exact witness checks and Proposition~\ref{prop:orientation} prove
the assertion for each configuration. Completeness of the list makes
every cell contribution in degree eight zero, giving
\eqref{eq:rank8}. The zero diagonal of \eqref{eq:rankfiltration}
and Hurewicz injectivity give \eqref{eq:localizedK}.

For the last step, localization zero means that each element of
\(K_8(\Z)\) is annihilated by a power of \(210\).
A nonzero such element would have a nonzero component of prime-power
order for one of \(2,3,5,7\), contrary to the cited arithmetic
vanishing. This implication does not require a separate rank
calculation or finite-generation argument.
\end{proof}

The initial program will thus check a concrete finite part of an
established proof while retaining an explicit list of its mathematical
premises. It will not certify that the 106 configurations exhaust all
orbits by reconstructing their enumeration. Removing that admission
would require a further completeness certificate for the classification.
Likewise, extending the program to higher \(K\)-groups requires new
mathematical input whenever the relevant chain groups or spectral
sequence diagonal no longer vanish.

\clearpage
\begin{thebibliography}{9}\small
\bibitem{DEKM}
M.~Dutour Sikirić, P.~Elbaz-Vincent, A.~Kupers and J.~Martinet,
\emph{Voronoi complexes in higher dimensions, cohomology of
\(\GL_N(\Z)\) for \(N\geq8\) and the triviality of \(K_8(\Z)\)}.
The theorem numbering used here follows the
\href{https://jamartin.perso.math.cnrs.fr/Publications/DEKM.pdf}{authors' manuscript};
see also \href{https://arxiv.org/abs/1910.11598}{arXiv:1910.11598}.

\bibitem{Data}
P.~Elbaz-Vincent, \emph{Voronoi complexes database},
\href{https://github.com/elbazvip/Voronoi-complexes-database}{public repository}.
The input inspected for this report is
\path{DATA_gap/ListSHV_n8_rnk9}, revision
\path{d4c7b40ad2c95850d16160b347953dbaa9241486}.
Its SHA-256 checksum is
\path{d0b58f52da26e64b20649dffe4e317efd31a9d6b0c5e4d5b1407dbf5d636359d}.
The implementation must retain and check the source bytes.

\bibitem{EVGS}
P.~Elbaz-Vincent, H.~Gangl and C.~Soulé,
\emph{Perfect forms, \(K\)-theory and the cohomology of modular groups},
Advances in Mathematics \textbf{245} (2013), 587--624.
\href{https://www.ihes.fr/~soule/documents/YAIMA4438.pdf}{Author PDF}.

\bibitem{Quillen}
D.~Quillen,
\emph{Finite generation of the groups \(K_i\) of rings of algebraic integers},
in \emph{Algebraic \(K\)-Theory I}, Lecture Notes in Mathematics
\textbf{341}, Springer, 1973, 179--198.

\bibitem{LS}
R.~Lee and R.~H.~Szczarba,
\emph{On the homology and cohomology of congruence subgroups},
Inventiones Mathematicae \textbf{33} (1976), 15--53.

\bibitem{Arlettaz}
D.~Arlettaz,
\emph{The Hurewicz homomorphism in algebraic \(K\)-theory},
Journal of Pure and Applied Algebra \textbf{71} (1991), 1--12.

\bibitem{Weibel}
C.~A.~Weibel, \emph{The \(K\)-book: An introduction to algebraic
\(K\)-theory}, Graduate Studies in Mathematics \textbf{145},
American Mathematical Society, 2013.
\href{https://sites.math.rutgers.edu/~weibel/Kbook/Kbook.VI.pdf}{Chapter VI}.
\end{thebibliography}
\end{document}
